\section{Waiting for the Bridegroom}
All ten virgins expect the bridegroom. All bring lamps, all become drowsy, and all wake at the cry. Their common expectation does not explain why five enter the wedding and five remain outside. Jesus places the difference in something prepared before the arrival: the wise have oil in their vessels. The foolish discover their need when there is no longer time to meet it. The concluding command to watch therefore requires more explanation than a command to remain awake. The wise have slept too.

The Fathers press the question into different parts of Christian life. John Chrysostom asks how someone who has undertaken the discipline of virginity can still neglect a neighbor. Augustine asks how even a person who feeds the hungry can lack charity. Gregory the Great examines the praise sought through good works, and, in his account of Emmaus, the hospitality through which disciples come to recognize Christ. These explanations concern acts, purposes, and the time in which conversion remains possible. They allow the parable to question both the person who has withheld help and the person whose help depends on being admired.

Wisdom approaches those who seek her; Paul consoles those who mourn their dead. Read beside the Gospel, these passages prevent readiness from becoming an exercise in anxious self-reliance. The one awaited is the risen Lord, and the ability to seek him begins with a gift. The Mass opens with a petition to be heard and ends by asking the Spirit to preserve those who have received. Its prayers concern people who need help throughout their response to God.

The first interpretation develops the urgency of mercy, especially Chrysostom's difficult explanation of the refusal to share oil. The second develops Augustine's distinction between a visible deed and the charity from which it proceeds. Their different accounts of the sellers remain visible: Chrysostom identifies them with the poor; Augustine and Gregory identify them with flatterers. Each explanation has consequences that cannot be recovered by treating the oil simply as a general symbol of goodness.

The map below groups the two lengths of the second reading and the two Communion antiphons as alternatives. No local choice is assumed. Scriptural Date and Location appears in Appendix~\ref{sec:date-location}; the passage commentary here concerns the texts' literary situations and the arguments made about them.
